% GENERATED FILE. Edit canonical lessons, then run npm run generate:books.
% canonical-source-hash: fnv1a64:43f089772c3716f6
% canonical-lessons: FA-C508-khoshnevisi, FA-C508-golduzi, FA-C508-qollab-bafi, FA-C508-kardasti, FA-C508-sofal

\chapter{Crafts}
\label{ch:fa-crafts}
\hlchaptermodality{508}

\begin{chapteropening}
\textbf{By the end of this chapter:} \emph{I can talk about crafts.}

Complete the last lesson of chapter 508: I can talk about crafts.
\end{chapteropening}

\section[\texorpdfstring{\fa{خوشنویسی} (khoshnevisi)}{خوشنویسی (khoshnevisi)}]{\fa{خوشنویسی} (khoshnevisi) — calligraphy}

\label{lesson:FA-C508-khoshnevisi}



\begin{quote}
Before the new one: say the Persian for to cook, then the Persian for a collection.
\end{quote}

\subsection*{You'll want to know: \fa{خوشنویسی}}
\textbf{\fa{خوشنویسی}} (\emph{khoshnevisi}) — \textquotedblleft{}calligraphy\textquotedblright{}.

\begin{quote}
\fa{او} \fa{سالها} \fa{خوشنویسی} \fa{یاد} \fa{گرفت}.
\end{quote}

\begin{quote}
\emph{u sâlhâ khoshnevisi yâd gereft.}
\end{quote}

\begin{quote}
He studied calligraphy for years.
\end{quote}

\emph{khosh} 'fine' + \emph{nevis} 'writing' (stem of \emph{neveshtan}) + \emph{-i}. A calligrapher is \emph{khoshnevis}.

\begin{cousinweb}
\emph{khoshnevisi} is built on native Persian parts: \emph{khosh}, 'fine', \emph{nevis}, the present stem of \emph{neveshtan}, which EtymDB traces to Middle Persian (EtymDB), and \emph{-i}: fine writing.
\end{cousinweb}

\begin{grammarlens}[title={What you've built}]
One more word for telling, explaining and describing at B1.
\end{grammarlens}

\subsection*{Your turn}
\begin{itemize}
\raggedright
  \item \emph{Say it:} \emph{khoshnevisi}
  \item \emph{Say it:} \emph{khoshnevisi}, once more
  \item \emph{Say it:} a sentence of your own with \emph{khoshnevisi}
  \item \emph{Recall:} say the Persian for to cook, then the Persian for a collection, then say \emph{khoshnevisi} again
\end{itemize}

\begin{tcolorbox}[breakable,colback=teal!4,colframe=teal!35!black,title={Before you move on}]
What does \fa{خوشنویسی} mean? (\textquotedblleft{}Calligraphy\textquotedblright{}.) Say it once more.
\end{tcolorbox}
\section[\texorpdfstring{\fa{گلدوزی} (golduzi)}{گلدوزی (golduzi)}]{\fa{گلدوزی} (golduzi) — embroidery}

\label{lesson:FA-C508-golduzi}



\begin{quote}
Before the new one: say the Persian for a collection, then the Persian for calligraphy.
\end{quote}

\subsection*{You'll want to know: \fa{گلدوزی}}
\textbf{\fa{گلدوزی}} (\emph{golduzi}) — \textquotedblleft{}embroidery\textquotedblright{}.

\begin{quote}
\fa{مادربزرگ} \fa{روی} \fa{دستمال} \fa{گلدوزی} \fa{میکند}.
\end{quote}

\begin{quote}
\emph{mâdarbozorg ru-ye dastmâl golduzi mikonad.}
\end{quote}

\begin{quote}
Grandmother embroiders handkerchiefs.
\end{quote}

\emph{gol} 'flower' + \emph{duz} 'sewing' (stem of \emph{dukhtan}) + \emph{-i}: flower-sewing.

\begin{cousinweb}
\emph{golduzi} is built on native Persian parts: \emph{gol}, 'flower', \emph{duz}, the present stem of \emph{dukhtan}, 'to sew', and \emph{-i}. EtymDB has no entry for the compound.
\end{cousinweb}

\begin{grammarlens}[title={What you've built}]
One more word for telling, explaining and describing at B1.
\end{grammarlens}

\subsection*{Your turn}
\begin{itemize}
\raggedright
  \item \emph{Say it:} \emph{golduzi}
  \item \emph{Say it:} \emph{golduzi}, once more
  \item \emph{Say it:} a sentence of your own with \emph{golduzi}
  \item \emph{Recall:} say the Persian for a collection, then the Persian for calligraphy, then say \emph{golduzi} again
\end{itemize}

\begin{tcolorbox}[breakable,colback=teal!4,colframe=teal!35!black,title={Before you move on}]
What does \fa{گلدوزی} mean? (\textquotedblleft{}Embroidery\textquotedblright{}.) Say it once more.
\end{tcolorbox}
\section[\texorpdfstring{\fa{قلاب} \fa{بافی} (qollâb bâfi)}{قلاب بافی (qollâb bâfi)}]{\fa{قلاب} \fa{بافی} (qollâb bâfi) — crochet}

\label{lesson:FA-C508-qollab-bafi}



\begin{quote}
Before the new one: say the Persian for calligraphy, then the Persian for embroidery.
\end{quote}

\subsection*{You'll want to know: \fa{قلاب} \fa{بافی}}
\textbf{\fa{قلاب} \fa{بافی}} (\emph{qollâb bâfi}) — \textquotedblleft{}crochet\textquotedblright{}.

\begin{quote}
\fa{خواهرم} \fa{برای} \fa{مادرم} \fa{یک} \fa{شال} \fa{قلاب} \fa{بافی} \fa{کرد}.
\end{quote}

\begin{quote}
\emph{khâharam barâ-ye mâdaram yek shâl qollâb bâfi kard.}
\end{quote}

\begin{quote}
My sister crocheted a scarf for my mother.
\end{quote}

\emph{qollâb} 'hook' + \emph{bâf} 'weaving' (stem of \emph{bâftan}) + \emph{-i}: hook-knitting.

\begin{cousinweb}
\emph{qollâb bâfi} is built in Persian on \emph{qollâb}, 'hook', a loan usually traced to Arabic, and \emph{bâf}, the stem of native \emph{bâftan}, 'to knit, to weave'. EtymDB has no entry for it.
\end{cousinweb}

\begin{grammarlens}[title={What you've built}]
One more word for telling, explaining and describing at B1.
\end{grammarlens}

\subsection*{Your turn}
\begin{itemize}
\raggedright
  \item \emph{Say it:} \emph{qollâb bâfi}
  \item \emph{Say it:} \emph{qollâb bâfi}, once more
  \item \emph{Say it:} a sentence of your own with \emph{qollâb bâfi}
  \item \emph{Recall:} say the Persian for calligraphy, then the Persian for embroidery, then say \emph{qollâb bâfi} again
\end{itemize}

\begin{tcolorbox}[breakable,colback=teal!4,colframe=teal!35!black,title={Before you move on}]
What does \fa{قلاب} \fa{بافی} mean? (\textquotedblleft{}Crochet\textquotedblright{}.) Say it once more.
\end{tcolorbox}
\section[\texorpdfstring{\fa{کاردستی} (kârdasti)}{کاردستی (kârdasti)}]{\fa{کاردستی} (kârdasti) — a handicraft, craft work}

\label{lesson:FA-C508-kardasti}



\begin{quote}
Before the new one: say the Persian for embroidery, then the Persian for crochet.
\end{quote}

\subsection*{You'll want to know: \fa{کاردستی}}
\textbf{\fa{کاردستی}} (\emph{kârdasti}) — \textquotedblleft{}a handicraft, craft work\textquotedblright{}.

\begin{quote}
\fa{بچه} \fa{ها} \fa{در} \fa{مدرسه} \fa{کاردستی} \fa{درست} \fa{کردند}.
\end{quote}

\begin{quote}
\emph{bachchehâ dar madrese kârdasti dorost kardand.}
\end{quote}

\begin{quote}
The children made crafts at school.
\end{quote}

\emph{kâr} 'work' + \emph{dast} 'hand' + \emph{-i}: handwork.

\begin{cousinweb}
\emph{kârdasti} is built on native Persian \emph{kâr}, 'work', and \emph{dast}, 'hand', both of which EtymDB traces to Middle Persian (EtymDB), with the ending \emph{-i}.
\end{cousinweb}

\begin{grammarlens}[title={What you've built}]
One more word for telling, explaining and describing at B1.
\end{grammarlens}

\subsection*{Your turn}
\begin{itemize}
\raggedright
  \item \emph{Say it:} \emph{kârdasti}
  \item \emph{Say it:} \emph{kârdasti}, once more
  \item \emph{Say it:} a sentence of your own with \emph{kârdasti}
  \item \emph{Recall:} say the Persian for embroidery, then the Persian for crochet, then say \emph{kârdasti} again
\end{itemize}

\begin{tcolorbox}[breakable,colback=teal!4,colframe=teal!35!black,title={Before you move on}]
What does \fa{کاردستی} mean? (\textquotedblleft{}A handicraft, craft work\textquotedblright{}.) Say it once more.
\end{tcolorbox}
\section[\texorpdfstring{\fa{سفال} (sofâl)}{سفال (sofâl)}]{\fa{سفال} (sofâl) — pottery, earthenware}

\label{lesson:FA-C508-sofal}



\begin{quote}
Before the new one: say the Persian for crochet, then the Persian for a handicraft.
\end{quote}

\subsection*{You'll want to know: \fa{سفال}}
\textbf{\fa{سفال}} (\emph{sofâl}) — \textquotedblleft{}pottery, earthenware\textquotedblright{}.

\begin{quote}
\fa{در} \fa{این} \fa{کارگاه} \fa{سفال} \fa{میسازند}.
\end{quote}

\begin{quote}
\emph{dar in kârgâh sofâl misâzand.}
\end{quote}

\begin{quote}
In this workshop they make pottery.
\end{quote}

\emph{sofâlgari} — the craft of pottery; \emph{sofâlgar} — a potter.

\begin{cousinweb}
\emph{sofâl} is the Persian word for earthenware, but EtymDB has no entry for it, and its origin is not settled.
\end{cousinweb}

\begin{grammarlens}[title={What you've built}]
One more word for telling, explaining and describing at B1.
\end{grammarlens}

\subsection*{Your turn}
\begin{itemize}
\raggedright
  \item \emph{Say it:} \emph{sofâl}
  \item \emph{Say it:} \emph{sofâl}, once more
  \item \emph{Say it:} a sentence of your own with \emph{sofâl}
  \item \emph{Recall:} say the Persian for crochet, then the Persian for a handicraft, then say \emph{sofâl} again
\end{itemize}

\begin{tcolorbox}[breakable,colback=teal!4,colframe=teal!35!black,title={Before you move on}]
What does \fa{سفال} mean? (\textquotedblleft{}Pottery, earthenware\textquotedblright{}.) Say it once more.
\end{tcolorbox}
